% 9-15(9쪽) — 조각마다 끊긴 함수 y=f(x) 의 그래프(선분 2개 · 곡선 2개 · 닫힌점·열린점). 스캔 화질이 낮아 TikZ 로 다시 그림.
% 원문에 있는 것만: 눈금 숫자 −2 −1 1 2 (x) · 2 1 −1 −2 (y) · 라벨 O, x, y, y=f(x) · 좌표 안내 파선.
% 한 칸은 원문처럼 가로·세로 같게(0.70cm) — 작게 잡으면 y축 −1 과 −2 라벨이 서로 겹친다.
\begin{tikzpicture}[x=0.70cm, y=0.70cm, line width=0.5pt]
  \draw[->, >=stealth, line width=0.7pt] (-2.7,0) -- (3.0,0) node[below=1pt] {$x$};
  \draw[->, >=stealth, line width=0.7pt] (0,-2.9) -- (0,2.9) node[left=1pt] {$y$};
  \node[below=2pt, font=\small] at (-0.3,0) {$\pt{O}$};
  % 좌표 안내 파선(원문)
  \draw[dashed, line width=0.35pt] (-1,1) -- (0,1);
  \draw[dashed, line width=0.35pt] (-1,1) -- (-1,-1) -- (1,-1);
  \draw[dashed, line width=0.35pt] (0.35,2) -- (2,2) -- (2,0);
  % 그래프
  \draw[line width=0.6pt] (-2,0) -- (-1,1);
  \draw[line width=0.6pt] (0,0) -- (1,-1);
  \draw[line width=0.6pt, smooth] plot coordinates {(-1,-1) (-0.87,-1.34) (-0.62,-1.66) (-0.36,-1.86) (0,-2)};
  \draw[line width=0.6pt, smooth] plot coordinates {(1,0) (1.10,0.54) (1.29,1.11) (1.55,1.56) (1.80,1.88) (2,2)};
  % 닫힌점
  \fill (-2,0) circle (2.2pt);
  \fill (-1,0) circle (2.2pt);
  \fill (0,-2) circle (2.2pt);
  \fill (1,-1) circle (2.2pt);
  \fill (2,2) circle (2.2pt);
  % 열린점
  \draw[fill=white, line width=0.5pt] (-1,1) circle (2.2pt);
  \draw[fill=white, line width=0.5pt] (-1,-1) circle (2.2pt);
  \draw[fill=white, line width=0.5pt] (0,0) circle (2.2pt);
  \draw[fill=white, line width=0.5pt] (1,0) circle (2.2pt);
  % 눈금 라벨(원문 자리 그대로: y 눈금은 모두 축 오른쪽)
  \node[right=2pt, font=\small] at (0,2) {$2$};
  \node[right=2pt, font=\small] at (0,1) {$1$};
  \node[below right=0pt and 3pt, font=\small] at (0,-1) {$-1$};
  \node[right=3pt, font=\small] at (0,-2) {$-2$};
  \node[below=2pt, font=\small] at (-2,0) {$-2$};
  \node[below=2pt, font=\small] at (-1.15,0) {$-1$};
  \node[below=2pt, font=\small] at (1.15,0) {$1$};
  \node[below=2pt, font=\small] at (2,0) {$2$};
  \node[right=1pt, font=\small] at (1.0,2.6) {$y=f(x)$};
\end{tikzpicture}
